\section{関連研究}

\subsection{単眼深度推定}
単眼深度推定は3Dモデリング\cite{4}，ロボット工学\cite{5}，自
動運転\cite{6}などに応用され，盛んに研究されている．これら
の手法はスケールが類似した大量のデータセットで学習・
テストを行うことで，学習時と同一の環境であれば距離推
定可能である．NDDepth\cite{1}やMIMDepth\cite{2}はKITTI\cite{7}
，NYUDepth-v2\cite{8}の2 種類のデータセットにおいて高い
精度を示した．しかし，学習した環境と異なる環境では，
深度から距離へと変換するためのスケールの推定が必要と
なる．したがってスケールの推定精度によっては精度が大
きく低下する可能性がある．
\subsection{距離補完}
入力にLiDAR 点群とカメラ画像を用いてLiDAR 点群を
補完し，密な距離画像を生成する距離補完手法も多く研究
されている\cite{9}\cite{10}．三次元地図点群を疑似的なLiDAR 点
群とみなすことで，距離画像の生成が可能である．しかし，
カメラ画像とLiDAR 点群は同時刻に取得されることが前提
の手法であり，画像内の物体とLiDAR 点群は一致している
ことが望まれる．そのため三次元地図上に存在しない車や
人などの動的障害物の距離推定においてPMOD-Net\cite{3}と
同等の精度を得るのは困難である．
\subsection{PMOD-Net}
PMOD-Net\cite{3}は地図から作成した疎な距離画像とカメ
ラ画像を入力として，Semantic Segmentation と距離推
定を行う世界初のニューラルネットワークである．
Semantic Segmentation と距離推定をマルチタスクに行
うことで三次元地図上に存在しない動的障害物を含めた三
次元障害物検出を可能にした．疎な距離画像は，自己位置
推定結果をもとに三次元地図点群を画像平面に投影するこ
とで取得する．自己位置はGPS やIMU，オドメトリ，カ
メラなどで取得した値を利用する．
\subsection{本研究の位置づけ}
PMOD-Netは，単眼カメラで障害物までの正確な距離を
計測できる利点がある．さらに実際に移動ロボットで障害
物回避にも成功しているため，本研究でもベース手法とし
て使用する．本研究は，単眼カメラのみでの代表的な距離
推定手法である単眼深度推定法とPMOD-Net\cite{3}を利用す
る．しかし，双方の相補的な関係に着目し，欠点を単一
ネットワーク内で同時に解決する新手法を開発したという
点で従来とは大きく異なる．

また本研究は，開発スピードが速い単眼深度推定分野に
おいて，深度推定法そのものを開発するものではない．地
図を利用することで，深度推定法によらずに距離推定を高
精度化する手法を開発している点も，従来とは大きく異な
る．さらに従来の学習を必要としないスケー
ル調整に加えて，学習ベースでスケールを調整するという新規性があり，より高精度な距離
推定を可能にしている．



